\section{Open questions}\label{sec:open}

The following five questions are the most consequential loose ends of this replication audit.
Each is grounded in a specific finding in \texttt{REPORT.md} or in the paper text; each has
concrete next steps that could be executed within a modest CPU budget (no HPC required for
any of them, except optionally item 5).

\subsection*{Q1 --- Is bare eq (2.5) actually stable at the published MCF7 fit?}

\textbf{Basis.} At MCF7's Table 1 constants, $k_3 k_5/(k_4 k_6) = 1.000374 > 1$; the post-repair
fixed point $(Y,Z) = (0,0)$ has $\det J = -145$, i.e.\ it is a saddle. Direct integration of
the printed bare eq (2.5) to 24 h yields $\langle Y\rangle = 3622$ vs $Y_{\max} = 300$ and
$\langle Z\rangle = 11313$ vs $Z_{\max} = 1000$. Table 2 caps are enforced by the SSA \emph{by
construction}, but they are absent from eq (2.5) as displayed. Paper does not discuss the
stability of the post-repair fixed point. MDA-MB-468 sits at ratio 0.987, marginally on the
stable side of the same bifurcation.

\textbf{Next steps.} (a) Contact the corresponding author (S.\ J.\ Chapman, Oxford) to ask
whether Y$_{\max}$/Z$_{\max}$ caps were applied silently in the Fig 4b integration and, if so,
which horizon. (b) If caps were applied, re-fit the six rate constants for MCF7 under the
constraint $k_3 k_5/(k_4 k_6) < 1$ and check whether Fig 4b fit quality degrades. (c) Publish
either an erratum clarifying the cap, or a corrected fit that lies strictly inside the stable
region.

\subsection*{Q2 --- Does the 6-parameter fit have predictive power beyond curve-fitting?}

\textbf{Basis.} Paper's eq (3.3) is minimised against Figure 4's foci/comet time series
(visually $\sim 7$ time points per cell line, 2 observables each, $\sim 14$ numbers) with 6
free rate constants. No cross-validation, no held-out data, no comparison to a 2-parameter
bi-exponential baseline. Section 5's discussion treats the fit as validation rather than as a
descriptive fit to be tested.

\textbf{Next steps.} (a) WebPlotDigitizer extraction of Figure 4a/b for both cell lines
($\sim 2$ h manual work). (b) Leave-one-out cross-validation: refit on 6 of 7 time points,
predict the seventh. (c) Head-to-head against $f(t) = A e^{-t/T_1} + B e^{-t/T_2}$ (2
parameters) on the same data with matched loss. (d) If the 6-parameter mechanistic model does
not substantially out-predict the 2-parameter descriptive baseline on held-out points, the
mechanistic interpretation of $k_1,\ldots,k_6$ is weakened.

\subsection*{Q3 --- How does this framework compare quantitatively against competing DSB-repair frameworks?}

\textbf{Basis.} Paper's Introduction \S1 cites prior mathematical models of DSB repair (two-lesion
kinetic, radiator/repair rate-equations, MEDRAS family) but does not run any of them on the
same MCF7/MDA-MB-468 datasets. No AIC/BIC across frameworks; no Bayes-factor comparison. This
is the single largest scientific-value gap.

\textbf{Next steps.} (a) Choose 1--2 competing frameworks with clean formulations (bi-exponential;
Cucinotta--Nikjoo two-lesion; a repair-priority-queue variant). (b) Fit each on the Fig 4
extracted data with matched loss. (c) Compute AIC, BIC, cross-validated NLL. (d) Report which
framework's inductive bias (single-Markov-site $+$ closures vs.\ two-lesion vs.\ empirical
bi-exponential) is best supported by the $\sim 14$-point data.

\subsection*{Q4 --- What is the quantitative mapping between SSA and ODE limits?}

\textbf{Basis.} Paper \S3 asserts the closures track SSA means ``well''; Fig 3 shows one overlay per
closure per cell line without a quantitative RMS or a convergence-in-$n$ table. Our local
\texttt{closure\_validation.py} gives RMS $Z \approx 56$ for MDA-MB-468 (SSA peak
$Z \approx 445$, $\sim 13\%$) and $\approx 104$ for MCF7 (SSA peak $Z \approx 800$,
$\sim 13\%$), noticeably looser than Fig 3's overlays suggest visually. Whether this is
tau-leap bias, true closure bias, or a difference in trajectory-count averaging is not
distinguished in the paper.

\textbf{Next steps.} (a) Run exact Gillespie for $n = 10, 100, 10^3, 10^4$ trajectories, both cell
lines, to isolate averaging noise from closure bias. (b) Compare exact-SSA mean at each $t$ vs
bare eq (2.5) and eq (2.8) with matched initial conditions. (c) Report RMS as a function of
$n$ and horizon. (d) If closure bias is $>5\%$ on the dynamic range of $\langle Z\rangle$, the
closures describe but do not quantitatively track the stochastic dynamics --- worth an
erratum-level correction.

\subsection*{Q5 --- Testable predictions at extremes: FLASH / ultra-low-dose / high-LET?}

\textbf{Basis.} Paper fits at 4 Gy X-ray dose (standard for MCF7/MDA-MB-468 in refs [12,17]).
Eq (4.3) Auger extension is exercised only at nominal specific-activity values in a small
range ($R = 0\ldots 8$). Paper does not test the framework at FLASH dose rate ($> 40$ Gy/s),
at ultra-low-dose ($< 0.01$ Gy), or at extreme-high-LET (Auger clusters, $\alpha$, HZE ions)
where the underlying assumptions (single-DSB dominance, Poisson DSB induction, density-independent
pATM pool) begin to break.

\textbf{Next steps.} (a) Push the SSA to a delta-function vs.\ boxcar DSB-induction regime
(FLASH-like). (b) Test a single-DSB-limit vs.\ a multi-DSB extension in which DSBs compete for
a finite pATM pool ($k_4$ becomes density-dependent). (c) Run the Auger case at cluster-DSB
parameter regimes where Poisson-induction assumption breaks. (d) Any excursion that shows a
\emph{qualitative} change in the AUC persistence measure defines a genuine testable prediction
worth taking to experiment. This is the highest-value open question for the framework's
scientific reach.
